\section{Methods}
\label{sec:methods}

\subsection{Waterline stacking to a bathymetric DEM}
\label{sec:dem}
Let $\{M_i\}$ be the Sentinel-1 binary water masks and $\{h_i\}$ their water levels,
sorted ascending by $h_i$. Pixels wet in the smallest mask (the always-wet core) are
assigned the lowest level $h_0$. A pixel that first becomes wet between consecutive
levels is a shoreline exposed in that increment and is assigned the mid-level
$(h_{i-1}+h_i)/2$. The resulting elevation field is closed over the union footprint,
interior gaps are filled by linear interpolation (\texttt{scipy.griddata}), and the
DEM is smoothed with a Gaussian filter ($\sigma=2$ px, $=20$~m). This adapts
\citet{Schwatke2020} to Sentinel-1 waterlines. The procedure is agnostic to the mask
source: any sequence of binary water masks with associated water levels stacks
identically, which we exploit to cross-validate the reconstruction with an independent
optical sensor (Section~\ref{sec:xsensor}). The DEM reconstructs only the
drawdown-\emph{exposed} band $[\,\text{floor},\text{max}\,]$; the permanently wet
deep zone is not observed.

\subsection{Scalar hypsometry and storage}
\label{sec:scalar}
As an alternative that needs only the area series, each SAR area is paired with the
nearest water level (within $\pm10$ days) and a power-law hypsometry is fitted,
\begin{equation}
  A(h) = a\,(h-h_0)^{b},
\label{eq:powerlaw}
\end{equation}
which is integrated to storage $V(h)=\int A\,\mathrm{d}h$. The DEM (Section~\ref{sec:dem})
gives the spatial bathymetry and change maps; the scalar fit gives a robust storage
curve where the observed level range is wide. Both see only the exposed band.

\subsection{Water-level-substitution ladder}
\label{sec:ladder}
To test a fully remote pipeline we keep the SAR area and the reconstruction fixed and
swap the water-level source between two independent tiers: the in-situ \emph{gauge}
and \emph{SWOT} altimetry. For each source we fit Eq.~\eqref{eq:powerlaw} and
integrate over the common observed band, then compare the resulting storage to each
other and to the field/survey truth. DAHITI is reported as a reprocessed-SWOT
cross-check (Section~\ref{sec:results}), not as an independent tier.

\subsection{Cross-sensor validation with PlanetScope optical masks}
\label{sec:xsensor}
Because the reconstruction (Section~\ref{sec:dem}) is agnostic to the mask source, we
independently rebuild the bathymetry from PlanetScope 3~m optical water masks and
compare it, pixel-wise and in AEV, to the Sentinel-1 reconstruction over the four
reservoirs PlanetScope covers (2024--2025). Water is delineated per scene by a
normalised difference water index,
$\mathrm{NDWI}=(\rho_\text{green}-\rho_\text{NIR})/(\rho_\text{green}+\rho_\text{NIR})$,
thresholded per scene, and the same gauge/SWOT water levels label the resulting
waterlines. This provides an independent, higher-resolution optical check and, for
\emph{Ancipa} and \emph{Pozzillo} (which lack a modern bathymetric survey), an
independent near-truth. Because optical detection carries error modes orthogonal to
SAR (surface algae and macrophytes, turbidity, sun glint), agreement between the two
reconstructions is strong evidence that neither is dominated by sensor-specific error.

\subsection{Capacity change: observable band versus total}
\label{sec:bandtotal}
Capacity change relative to the design curve is quantified two ways: \emph{band}
(relative to the DEM floor, i.e.\ the SAR-observable exposed zone, the fair
comparison against survey truth) and \emph{total} (the full design curve, including
the deep zone below the floor that SAR cannot see). Because sedimentation concentrates
in the deep zone, the SAR band change is a validated \emph{lower bound} on total loss.

\subsection{Uncertainty and bias}
\label{sec:uncertainty}
The one field-validated reservoir (Garcia) yields a normalised, elevation-dependent
bias model $\mathrm{bias}(f)$, $f=(z-\text{floor})/(\text{max}-\text{floor})$, which
is applied to the other reservoirs and propagated through a coherent Monte-Carlo to
put a confidence interval on the storage estimate. Because per-pixel error averages
out at reservoir scale, the dominant uncertainty is the systematic bias-transfer term.

\subsection{Interactive tool}
\label{sec:tool}
The pipeline is released as an interactive application (Streamlit + Earth-Engine
Python) that, for a selected reservoir, displays the reconstructed 2D/3D bathymetry,
the AEV curves against design and survey references, the multi-period change map, and
a downloadable DEM.
